% appC-online.tex -- full deferred proofs of Section 4.
\section{Deferred proofs of \Cref{sec:online}}\label{app:online}

\subsection{The lower bound (\Cref{lem:online-lower})}\label{app:online-lower-full}

Fix a signing $\xi:[k]\to\R$ with $\sum_c\xi(c)=0$ and apply it coordinatewise to an online $k$-coloring: at each step the
sign assigned to the arrival is $\xi(\chi(t))$, a function of the (online) color, hence itself online. The resulting signed
sum at coordinate $i$ is $\sum_c\xi(c)x_{i,c}$. For the even split ($\xi=\pm1$ on equal halves), $\sum_c\xi(c)x_{i,c}$ is a
difference of two sums of $k/2$ loads, each within $\rangeop_i$ of the others, so $|\sum_c\xi(c)x_{i,c}|\le(k/2)\rangeop_i$.
This is a valid online $\pm1$-balancing of the same arrivals (a uniform $d$-sparse stream), so the two-color online lower
bound $\Om(\loglog n)$ of \cite{AT2025} gives $(k/2)\max_{t,i}\rangeop_i\ge\Om(\loglog n)$, i.e.\ $\ondisc_k\ge\Om(\loglog n/k)$.

For odd $k=2a+1$, take $\xi\equiv1$ on a part $P_+$ of size $a+1$ and $\xi\equiv-\tfrac{a+1}{a}$ on $P_-$ of size $a$ (zero-sum).
After rescaling, the induced online process uses step sizes in $\{+1,-\rho\}$ with $\rho=a/(a+1)\in[\tfrac12,1]$. The
two-color lower bound of \cite{AT2025} is driven by a recurrence on the spread that, at each of its $\Th(\loglog n)$ stages,
grows by a constant amount independent of the step size within the band $[\tfrac12,1]$; the only effect of $\rho\ne1$ is a
constant factor $\le2$ in the stage count. Hence $\Om(\loglog n)$ persists for the projected process and
$\ondisc_k\ge\Om(\loglog n/k)$ for odd $k$ as well. \qed

\subsection{The upper bound (\Cref{lem:online-upper})}\label{app:online-upper-full}

Write $\Hscale=\loglog n$. Let $\widetilde x_{i,c}(t)$ be the loads under the pure-seed coloring ($\chi\equiv\widetilde\chi$)
and $\widetilde R_i(t)$ the seed-only range. A coordinate--arrival pair $(i,t)$ is a \emph{modified incidence} if $i\in S_t$
and the algorithm's color differs from $\widetilde\chi_t$; by construction the color differs only when the repair branch
fires, which requires a coordinate $i^*\in S_t\cap E_{t-1}$. Since range is $2$-Lipschitz in $\ell_\infty$,
$|R_i(t)-\widetilde R_i(t)|\le 2\,\|x_i(t)-\widetilde x_i(t)\|_\infty\le 2\cdot(\#\text{ modified incidences at }i)$.

\paragraph{Part 1: seed tail.}
Fix $i$ and colors $a\ne b$, and let $\widetilde D_{i,a,b}(t)=\sum_{s\le t}\1[i\in S_s]\big(\1[\widetilde\chi_s=a]-\1[\widetilde\chi_s=b]\big)$,
a martingale with increments in $\{-1,0,1\}$ and predictable variance $\langle\widetilde D\rangle_T\le\sum_s\Pr(i\in S_s)\Pr(\widetilde\chi_s\in\{a,b\})=2Td/(nk)=:V$.
By the Freedman/Bernstein maximal inequality,
$\Pr(\max_t\widetilde D_{i,a,b}\ge 32\Hscale)\le\exp(-\tfrac{(32\Hscale)^2/2}{V+32\Hscale/3})$. The seed range exceeds
$32\Hscale$ only if some pair $(a,b)$ has $\widetilde D_{i,a,b}\ge 32\Hscale$, so a union over $\le k^2$ pairs gives
$\Pr(\max_t\widetilde R_i\ge 32\Hscale)\le k^2\exp(-\tfrac{512\Hscale^2}{V+\Oh(\Hscale)})$. In the worst case $k=2$,
$d=\Hscale^2/\log\Hscale$: $V=\Hscale^2/\log\Hscale$ dominates the $\Oh(\Hscale)$ term, the exponent is $\approx512\log\Hscale$,
and the bound is $\Hscale^{-512+o(1)}\le d^{-28}$ (since $d^{-28}=\Hscale^{-56+o(1)}$): a factor-$\approx 9$ margin in the
exponent. Increasing $k$ lowers $V$ and only helps. Thus
\begin{equation}\label{eq:seed-tail}
  \Pr(\max_t\widetilde R_i\ge 32\Hscale)\ \le\ d^{-28}.
\end{equation}

\paragraph{Parts 2--3: exceptional set and decay, by reduction to the two-color analysis.}
The remaining two steps are exactly the category and level-set analysis of \cite{AT2025}, which we invoke as a black box.
Their analysis applies to any online algorithm built from a random seed coloring together with a deterministic repair rule
satisfying:
\begin{itemize}
  \item[(i)] the repair never raises the range of the coordinate it repairs; and
  \item[(ii)] the assigned color differs from the seed only on arrivals whose support contains a coordinate already at range
  $\ge 64\Hscale$.
\end{itemize}
Under (i)--(ii) and the seed tail \eqref{eq:seed-tail}, \cite[exceptional-set lemma]{AT2025} gives $|E_T|\le n d^{-5}$ \whp,
and \cite[level-set lemma]{AT2025} gives, for $M(t,r)=\{i:\max_{s\le t}\rangeop_i(s)\ge 65\Hscale+3r\}$,
$|M(T,r)|\le n(\log n)^{-2\cdot 2^r}$, which is empty by $r=\lceil 2\Hscale\rceil$, whence $\max_{t,i}\rangeop_i<80\Hscale$.
Both lemmas use only the independence of future supports, the seed tail, and (i)--(ii); the number of colors enters nowhere
(the category depth is indexed by arrivals, not colors, and the level-set recurrence counts same-arrival exceptional
helpers).

It remains to check that our $k$-color algorithm satisfies (i)--(ii). For (i): a minimum-load color of $i^*$ is, by
definition, a coordinate of $x_{i^*}$ attaining the minimum, and incrementing a minimum coordinate cannot increase
$\max_c x_{i^*,c}-\min_c x_{i^*,c}$ for any $k$ (when several colors tie for the minimum, any one works). For (ii): the
assigned color equals the seed $\widetilde\chi_t$ unless the repair branch fires, and that branch fires only when
$S_t\cap E_{t-1}\ne\varnothing$, i.e.\ when the support contains a coordinate already at range $\ge 64\Hscale$. Both
properties are statements about the algorithm's control flow and the range function, with no dependence on $k$. Hence the
two lemmas of \cite{AT2025} apply verbatim, and $\max_{t,i}\rangeop_i<80\Hscale$ \whp. \qed

\begin{remark}
The reduction is the natural one: the only place the number of colors could matter is the repair, and a minimum-load repair
is range-non-increasing for every $k$. The modified-incidence count that drives \cite{AT2025}'s category recursion charges
each deviation of the actual coloring from the seed to a co-present exceptional coordinate, and this charging is unchanged by
$k$ because it depends only on the binary event ``actual color $\ne$ seed'', not on which color was used.
\end{remark}
